\documentclass[../main.tex]{subfiles}

\begin{document}
\section{Introduction}

Nowadays, geometric modelers are everywhere to be found, for example, in
archaeology, architecture and mechanical engineering, just to name a few
domains. Their importance in design and conception have grown to make them
essentials. These softwares rely on geometric modeling kernels to maintain the
designs and allow precise and complex works.

Topological changes (or \emph{events}) tracking is a key functionality for these
kernels such as (\textit{ACIS}, …). Such a functionality shows its importance
during the construction of a model when the applied operations modify its
geometry and topology. Considering how models tend to grow during their
conception, the tracking may become slower and coherence issues may arise.



\subsection*{motivation}
\begin{itemize}
    \item Modelers API (events logging) \cite{baba2007generic}
    \item Persistent naming \cite{cardot2019persistent} \cite{jauhar2020assembly} and model re-evaluation/re-construction \cite{guo2022complexgen} \cite{lambourne2022reconstructing}
    \item parallelism (automatique ou non) \cite{gargallo2017subdividing} \cite{bourquat2020transparent} \cite{mlakar2018alsub} \cite{kalatzantonakis2020cooperative}
    \item …
\end{itemize}
\subsection*{Verrou}
\begin{itemize}
    \item méthode dynamique (coûteuse temps/place)
    \item méthode dynamique partielle (nommage, par opération)
\end{itemize}
%
\subsection*{Approche proposée}

Our approach relies on the static analysis of graph transformation rules to detect events before their applications on an object.

\section{Related works}

\end{document}

%%% Local Variables:
%%% mode: latex
%%% TeX-master: "../main"
%%% End:
